\section{Introduction}\label{sec:intro}

Production retrieval-augmented question-answering systems are
increasingly built on top of knowledge graphs (KGs). For domains where
facts are inherently multi-relational — public-sector regulations,
clinical guidelines, organizational policies — the canonical argument
for a KG backbone runs as follows. Natural-language documents lose
the schema-driven structure that the regulation \emph{has} as a fact
base. A binary triple store flattens N-ary relations and forces
either lossy decomposition or awkward reification. A
hyper-relational KG (e.g., TypeDB v3) preserves the original role-
player structure faithfully, and a ReAct-style agent can in principle
exploit that structure via formal queries (TypeQL, Cypher) for
multi-hop reasoning.

\paragraph{Scope.} Our setting is a regulatory KG-RAG deployment in
which the entire structured fact base fits within a modern LLM's
32K-token context window. We argue this scope is unusually broad in
practice: regulations, organizational policies, clinical guidelines,
benefits schedules, and many similar enterprise corpora have fact
bases of low single-digit megabytes once serialized, well below
state-of-the-art context limits. The claims in this paper are
explicitly scoped to such fact-bounded domains; multi-million-fact
KGs (knowledge bases of the open Web, biomedical corpora, etc.)
require retrieval and are out of scope for the surface-form ablation.

This paper reports a counterintuitive finding from a production
deployment of exactly such a fact-bounded system. Over the Bank of
Korea Compensation Regulations — a 200-page Korean public-sector
document with eight N-ary fact tables and twelve supplementary
clauses — we evaluate four LLM families against eight backends (TypeDB
hyper-relational agent, Neo4j property-graph agent, production
chunked-context retriever, closed-book LLM, and four
\emph{in-context surface-form} variants in which the entire KG is
serialized into the prompt as TypeQL, Cypher, JSON-LD, or natural
language). The benchmark is 100 questions covering ten categories,
including 42 multi-key questions that stress the supposed strength
of the hyper-relational KG.

The headline result is that the production KG agents
\emph{underperform} every in-context surface form by 8--26
percentage points (paired McNemar $p<0.05$ in 5 of 8 KG-vs-context
comparisons). Three observations make this counterintuitive:
(i) the KG agents have access to \emph{strictly more facts} than
the in-context dumps; (ii) the gap is largest exactly in the
multi-key categories that the KG was designed to serve; and
(iii) the closed-book baseline is at 9--13\%, ruling out
pre-training memorization as the explanation.

The surface-form ablation localizes the bottleneck. Four findings
result.

\begin{enumerate}
\item \textbf{The query-generation stage, not KG expressivity, is
      the dominant accuracy determinant.} An in-context dump with
      the same hyper-relational structure as the KG store reaches
      83--95\% accuracy across LLMs, while the production agent
      built on top of the KG reaches 56--73\%. The 8--26 percentage
      points lost between the two are absorbed by the agent's
      query-authoring, query-execution, and result-verbalization
      stages.
\item \textbf{Natural-language summary Pareto-dominates structured
      surface forms.} Across LLMs, NL achieves a mean accuracy of
      89.8\% at 14K characters, beating the structured tiers
      (TQL/Cypher/JSON-LD) at 83--85\% and 26--32K characters.
\item \textbf{Llama-4-Scout exhibits a 21-percentage-point form
      sensitivity that other LLMs do not.} The asymmetry traces to
      query-language familiarity in pre-training, with single-fact
      lookups surviving unfamiliar syntax but cross-reference tasks
      collapsing — a finding directly testable via guide ablation.
\item \textbf{In-context KG paradoxically underperforms NL on
      closed-world boundary inference.} The structural strength of
      the KG (schema-driven clarity, integrity rules) is stripped
      away when serialized; the resulting raw dump leaves the LLM
      to enumerate-and-disprove rather than read an explicit
      boundary clause.
\end{enumerate}

We argue for a reframing of KG infrastructure value: the KG should
be treated as the \emph{operational backbone} of the system —
schema enforcement, multi-system integration, latent expressivity
for stronger query agents — rather than as the LLM-facing surface
form. As a practical recommendation, we propose adding a
KG-to-natural-language conversion layer between the KG store and
the LLM, augmented with explicit boundary annotations to address
the closed-world inference gap. Section~\ref{sec:discussion}
develops these implications.

\subsection*{Contributions}

\begin{itemize}
\item We report the first production-deployment evaluation of
      TypeDB v3 hyper-relational and Neo4j property-graph KGs
      against a four-tier in-context surface-form ablation
      (Section~\ref{sec:results}).
\item We demonstrate that the query-generation stage is the
      dominant accuracy bottleneck in current KG-RAG agents,
      not KG expressivity (Section~\ref{sec:results}).
\item We identify and decompose two distinct surface-form
      mechanisms — query-language familiarity and closed-world
      boundary inference — and show that they require distinct
      remedies (Section~\ref{sec:mechanism}).
\item We articulate the \emph{KG semantic paradox} and recommend
      a KG-to-NL conversion layer as the practical next step for
      KG-LLM pipelines (Section~\ref{sec:discussion}).
\end{itemize}
